\begin{finding}[\fid{F25}: the commutant defect, coherent extensions, and a regular basis]\label{f:r3deformation}
Use \fid{F21} over $R=\Q[v,v^{-1}]_{(v-1)}$ and write
$\Delta=z_1-z_2$. Then
\[
 \End_A(M)=M_{20}(R)\times M_4(R)\times M_{16}(R)\times M_4(R).
\]
This is free of rank $688$. Its base-change map to the full special
commutant
\[
 \End_{A_1}(M_1)=M_{20}(k)\times M_4(k)\times M_{20}(k)
\]
is injective with cokernel of dimension $128$; the latter commutant has
dimension $816$. Moreover
\[
 \operatorname{Ext}^1_C(L_1,L_2)
 =\operatorname{Ext}^1_C(L_2,L_1)=R/(\Delta),\qquad
 \operatorname{Ext}^1_C(L_i,L_i)=0,
\]
and $\operatorname{Ext}^1_A(M,M)=(R/(\Delta))^{128}$, of $R$-length $256$.

There is also a regular center basis $1,y$, with
\begin{equation}\label{eq:r3basis}
 y=D^2(z_1-z),\qquad
 y^2=\lambda y,\qquad
 \lambda=(v-v^{-1})^2(v^2+1+v^{-2}).
\end{equation}
Together with the two scalar-block units, $E_{ij},yE_{ij}$ form a basis
of the actual rank-ten algebra over $R_0$, with structure constants
among $0,1,\lambda$. This elementary basis is not asserted to be VIII's
designated canonical basis.
\end{finding}
\begin{proof}
Central blocks separate the two scalar sectors and the matrix sector.
Morita equivalence identifies the last sector's multiplicity module
with $L_1^{16}\oplus L_2^4$. A homomorphism $L_1\to L_2$ is determined
by $r\in R$ satisfying $(z_2-z_1)r=0$. Since $\Delta$ is a nonzero
element of the domain $R$, this Hom group is zero, as is the reverse
one; $\End_C(L_i)=R$. This gives the asserted commutant. At $v=1$ the
two branch modules become the same simple residue module. Their
multiplicity spaces merge into $k^{20}$. The map on these factors is
the block-diagonal embedding $M_{16}(k)\times M_4(k)\hookrightarrow
M_{20}(k)$; the missing off-diagonal blocks each have dimension $16\cdot4$.

For the extension calculation put $f_i=z-z_i$. The following is an
exact free resolution of $L_1$:
\[
 \cdots\xrightarrow{f_1}C\xrightarrow{f_2}C
             \xrightarrow{f_1}C\longrightarrow L_1\longrightarrow0.
\]
To check exactness, lift to the domain $R[z]$ and cancel $f_1$ or $f_2$
from divisibility by $f_1f_2$. Applying $\Hom_C(-,L_2)$ gives
$R\xrightarrow{-\Delta}R\xrightarrow0R\xrightarrow{-\Delta}R\to\cdots$,
so degree one is $R/(\Delta)$. Exchanging indices gives the reverse
group. Applying $\Hom_C(-,L_1)$ instead begins
$R\xrightarrow0R\xrightarrow\Delta R$ and has zero degree one.
Morita equivalence and the multiplicities in \eqref{eq:r3tensor} give
$16\cdot4+4\cdot16=128$ copies; scalar blocks contribute no extensions
of their free $R$-modules. Since $\Delta$ is a unit times $(v-1)^2$,
each crossed quotient has length two. For the general $C_m$ the same
resolution gives crossed extension group $R/(t^m)$, retaining precisely
the collision order omitted by the constant \'etale data.

Finally $D$ is a unit and the affine change $z\mapsto y$ is invertible.
The two evaluations of $y$ are $0$ and $D^2(z_1-z_2)$. Their product
relation yields $y^2=D^2\Delta y$. Since
\[
 D^2\Delta=(D-1)(D-4),\quad
 D-4=(v-v^{-1})^2,\quad D-1=v^2+1+v^{-2},
\]
this is \eqref{eq:r3basis}. Matrix-unit multiplication proves the
structure-constant assertion. At $v=1$, $\lambda=0$ and the four
$yE_{ij}$ span the radical.
\end{proof}
IV \S12.3 already records the numbers $688$ and $816$. Here the free
rank-$688$ commutant family is kept distinct from the full
special-fiber commutant, and their comparison map is explicit.

The scalar in \eqref{eq:r3basis} has a precise ordinary geometric trace:
with $Q=v^2$,
\[
 \lambda=Q^{-2}(Q-1)^2(1+Q+Q^2)
       =Q^{-2}\#(\mathbb G_m^2\times\PP^2)(\F_Q).
\]
It is the compact-support alternating trace with Tate twist $(2)$.
The remaining factor $v^{-2}(1+v^2+v^4)$ is the centered Tate polynomial
of $\PP^2$. This mixed-cohomology trace interpretation of one scalar
does not supply functorial multiplication or associativity constraints,
and it does not assert VIII's consecutive-$v$ unimodality convention.
Indeed $\lambda=v^4-v^2-v^{-2}+v^{-4}$ has negative Laurent coefficients.


